\chapter{Cyclotomic values and the polygamma bridge}
\label{ch:02-cyclotomic}

Roots of unity turn additive Fourier coordinates into arithmetic character
coordinates. The same finite transform carries Hurwitz zeta, positive-order
polygamma and Clausen values. It explains which half of a polylogarithm is
elementary at a given weight and which character values remain to be studied.

\section{The two-way exact DFT}\label{bridge:sec:dft}

We use the Clausen convention
\[
\Cl_{2m}(\theta)=\Im\Li_{2m}(e^{i\theta}),\qquad
\Cl_{2m+1}(\theta)=\Re\Li_{2m+1}(e^{i\theta}).
\]
The sine series at odd weight is a different function; its rational-angle
values reduce to powers of $\pi$ through Bernoulli polynomials.

\begin{law}[forward direction; classical]\label{bridge:law:forward}
For every $n\ge2$, every $q\ge1$, and $\gcd(p,q)=1$,
\[
\Li_n\!\bigl(e^{2\pi ip/q}\bigr)
=\frac{(-1)^n}{q^n\,(n-1)!}\sum_{k=1}^{q-1}e^{2\pi ikp/q}\,
\psi^{(n-1)}\!\Bigl(\frac{k}{q}\Bigr)\;+\;\frac{\zeta(n)}{q^n}.
\]
\end{law}

This is identity 10.08.03.0007.01 of the \texttt{functions.wolfram.com} \emph{PolyLog}
reference \cite{bridge:WolframFunctions}; it is provable by rearranging the absolutely convergent Hurwitz-zeta series, and
it is the workhorse of the project's cyclotomic reducer (\texttt{cycloExact}): a finite
polygamma sum, exact and fast, resolving $\Li_n$ at \emph{every} root of unity with no
PSLQ, no squarefree-radical basis, and no nested-radical limitation.

\begin{result}[inverse direction]\label{bridge:res:inverse}
Inverting the discrete Fourier transform over $\mathbb{Z}/q$ gives, with
$\zeta_q=e^{2\pi i/q}$,
\[
\psi^{(n-1)}\!\Bigl(\frac{k}{q}\Bigr)
=(-1)^n\,(n-1)!\;q^{\,n-1}\sum_{p=0}^{q-1}\zeta_q^{-pk}\,\Li_n\!\bigl(\zeta_q^{\,p}\bigr).
\]
\dps{verified to $40$ digits at weights $2,3$ across $q\le21$; re-verified by mpmath at
$70$ dps for $n-1=1,2,3$ over the full store grid}
\end{result}

The bridge is therefore an \emph{exact finite Fourier transform, in both directions}.
Three non-obvious consequences are recorded.

\begin{enumerate}
\item \textbf{DFT-conjugacy of the relation lattices.} The distribution relation
$\sum_{j}\Li_n(z\,\zeta_m^{\,j})=m^{1-n}\Li_n(z^m)$ Fourier-transforms into the polygamma
multiplication theorem, and polylogarithm inversion/reflection corresponds to the symmetry
$k\leftrightarrow q-k$ of the $\psi$-grid. This is \emph{why} the two relation lattices ---
computed independently by the project's two reducers --- carry identical dimension counts
$\varphi(q)/2$ per denominator grid.
\item \textbf{The parity split, and Catalan as the $q=4$ odd component.} The imaginary parts are the odd Fourier sector under $k\mapsto q-k$.
The usual even-weight Clausen values belong to this sector, while odd-weight
Clausen values are real parts in the even sector. In particular
\[
\Catalan=\operatorname{Im}\Li_2(i)=\frac{\psi^{(1)}(1/4)-\psi^{(1)}(3/4)}{16},
\]
which identifies the well-known ``arguments $1/4$ and $3/4$'' pairing (Math.SE 893486) as
the $q=4$ odd component of the parity split.
\item \textbf{Certificate transport.} Cyclotomic-polylogarithm proof certificates can in
principle be transported to and from $\psi$-side lattice certificates through the exact
$\mathbb{Q}(\zeta_q)$ DFT matrix. The $\psi$-side elementary atoms ($\pi^2\csc^2$ values)
are far lighter than the sine-product algebraic objects that constrain the
$\Gamma$-certificate sweeps on large grids, so the transport is a proposed escape route
from that bottleneck (not yet implemented).
\end{enumerate}

\section{Clausen-type laws for polygamma values at rational arguments}\label{bridge:sec:laws}

MathWorld's \emph{Polygamma Function} page lists fifteen closed forms in its
special-cases block, eqs.~(22)--(36) ($\psi^{(1)}$, $\psi^{(2)}$, $\psi^{(3)}$ at $p/q$
with $q\in\{2,3,4,6\}$) in Catalan/Clausen/$\beta$ atoms, several of them traditionally
presented as isolated curiosities. All of them ---
and their extension to \emph{every} rational argument --- are instances of three
sine/cosine-weighted Clausen-DFT laws, obtained by splitting the inverse DFT of
Result~\ref{bridge:res:inverse} into its elementary (reflection) half and its Clausen half.

\begin{law}[trigamma; Lewin 1991]\label{bridge:law:psi1}
For $0<p<q$, $\gcd(p,q)=1$,
\[
\psi^{(1)}\!\Bigl(\frac{p}{q}\Bigr)
=\frac{\pi^2}{2\sin^2(\pi p/q)}
+2q\!\!\sum_{j=1}^{\lfloor(q-1)/2\rfloor}\!\!\sin\Bigl(\frac{2\pi jp}{q}\Bigr)\,
\Cl_2\Bigl(\frac{2\pi j}{q}\Bigr).
\]
\dps{verified at $q=3,4,5,8,12$ to $10^{-46}$ or better, then over the full store grid to
$10^{-50}$ in-kernel and at $70$ dps in mpmath}
\end{law}

This law is classical --- MathWorld's \emph{Trigamma Function} eq.~(14) cites Lewin
(1991) \cite{bridge:Lewin} --- and was independently re-derived here from the DFT bridge before the attribution
was spotted; the project's records explicitly re-frame the finding as \emph{placement},
not novelty. MathWorld's scattered coefficients are its small-$q$ shadows: the $3\sqrt3$
of $q=3$ is $2q\sin(2\pi/3)$, the $8\Catalan$ of $q=4$ is $2q\,\Cl_2(\pi/2)$.

\begin{law}[tetragamma]\label{bridge:law:psi2}
\[
\psi^{(2)}\!\Bigl(\frac{p}{q}\Bigr)
=-\pi^3\cot\Bigl(\frac{\pi p}{q}\Bigr)\csc^2\Bigl(\frac{\pi p}{q}\Bigr)
-2q^2\Biggl(\zeta(3)
+2\!\!\sum_{j=1}^{\lfloor(q-1)/2\rfloor}\!\!\cos\Bigl(\frac{2\pi jp}{q}\Bigr)
\Cl_3\Bigl(\frac{2\pi j}{q}\Bigr)+\varepsilon_q\Biggr),
\]
where $\varepsilon_q=(-1)^p\,\Cl_3(\pi)$ for even $q$, $\varepsilon_q=0$ for odd $q$,
and $\Cl_3(\pi)=-\tfrac34\zeta(3)$.
\dps{verified over the full store grid to $10^{-50}$ in-kernel and at $70$ dps in mpmath}
\end{law}

\begin{law}[pentagamma]\label{bridge:law:psi3}
\[
\psi^{(3)}\!\Bigl(\frac{p}{q}\Bigr)
=\pi^4\Bigl(2\cot^2\Bigl(\frac{\pi p}{q}\Bigr)\csc^2\Bigl(\frac{\pi p}{q}\Bigr)
+\csc^4\Bigl(\frac{\pi p}{q}\Bigr)\Bigr)
+12\,q^3\!\!\sum_{j=1}^{\lfloor(q-1)/2\rfloor}\!\!\sin\Bigl(\frac{2\pi jp}{q}\Bigr)\,
\Cl_4\Bigl(\frac{2\pi j}{q}\Bigr),
\]
equivalently, in difference form,
$\psi^{(3)}(p/q)-\psi^{(3)}(1-p/q)=24\,q^3\sum_j\sin(2\pi jp/q)\,\Cl_4(2\pi j/q)$.
\dps{same verification battery}
\end{law}

The $\psi^{(3)}$ law is the Lewin law one weight up, by the same derivation; no literature
citation has been located, and the project records it as ``presumably in Lewin's orbit.''
MathWorld's $162\sqrt3$ and $768\beta(4)$ coefficients are $12q^3\sin(2\pi j/q)$-instances
of the one-sided law ($24q^3$ in the difference form).

\begin{result}[new instances beyond the MathWorld table]\label{bridge:res:newinstances}
Among the instances at $q\in\{5,8,12\}$:
\[
\psi^{(1)}\!\bigl(\tfrac15\bigr)
=\frac{\pi^2}{2\sin^2(\pi/5)}
+10\Bigl[\sin\tfrac{2\pi}{5}\,\Cl_2\!\bigl(\tfrac{2\pi}{5}\bigr)
+\sin\tfrac{4\pi}{5}\,\Cl_2\!\bigl(\tfrac{4\pi}{5}\bigr)\Bigr],
\]
\[
\psi^{(1)}\!\bigl(\tfrac1{12}\bigr)
=40\,\Catalan+\frac{4\pi^2}{(\sqrt3-1)^2}+30\sqrt3\,\Cl_2\!\bigl(\tfrac{2\pi}{3}\bigr),
\]
the latter after duplication/triplication canonicalization collapses the raw five-atom DFT
form onto the canonical atom set. Also, MathWorld's eq.~(27) simplifies: from
$\Cl_3(2\pi/3)=-\tfrac49\zeta(3)$,
\[
\psi^{(2)}\!\bigl(\tfrac13\bigr)=-\frac{4\pi^3}{3\sqrt3}-26\,\zeta(3).
\]
\end{result}

\begin{result}[the polygamma--Clausen store]\label{bridge:res:store}
The corpus now carries \texttt{identities/polygamma-clausen-values.wl}: $54$ entries
$\psi^{(m)}(p/q)$ for $m\in\{1,2,3\}$, $q\in\{3,4,5,6,8,12\}$, all $p$ coprime to $q$
($18$ arguments $\times\,3$ orders), in canonicalized Clausen atoms. The residual atom
basis across the store is
\begin{align*}
&\{\Catalan,\ \beta(4),\ \zeta(3)\}\ \cup\
\Cl_2\bigl\{\tfrac{2\pi}{3},\tfrac{2\pi}{5},\tfrac{4\pi}{5},\tfrac{\pi}{4}\bigr\}\\
&\qquad\cup\
\Cl_3\bigl\{\tfrac{2\pi}{5},\tfrac{\pi}{4},\tfrac{\pi}{6}\bigr\}\ \cup\
\Cl_4\bigl\{\tfrac{2\pi}{3},\tfrac{2\pi}{5},\tfrac{4\pi}{5},\tfrac{\pi}{4}\bigr\}.
\end{align*}
Every entry is verified in-kernel to $|{\rm lhs}-{\rm rhs}|<10^{-50}$ at $60$ digits
(including a disk round-trip) and independently by an mpmath script at $70$ dps
(tolerance $10^{-55}$) that re-derives the raw DFT bridge, the three laws, all $17$
Clausen canonicalization relations, and the $L$-function bridges of
Section~\ref{bridge:sec:collapse} from mpmath primitives, with positive and negative PSLQ
controls.
\end{result}

\section{The boundary of the odd-weight Clausen collapse}\label{bridge:sec:collapse}

A folklore rule says that odd-weight Clausen values at rational angles collapse to
rational multiples of zeta values ($\Cl_3(2\pi/3)=-\tfrac49\zeta(3)$ being the standard
example). The store emission \emph{sharpened the rule's scope}:

\begin{proposition}[Crystallographic-angle reductions]
At denominators $q\in\{1,2,3,4,6\}$, the odd Clausen value
$\Cl_3(2\pi j/q)=\Re\Li_3(e^{2\pi i j/q})$ reduces to a rational multiple
of $\zeta(3)$. At $q=5,8,12$ the character decomposition instead gives
\begin{gather*}
\Cl_3(2\pi/5)=-\frac6{25}\zeta(3)+\frac{\sqrt5}{4}L(3,\chi_5),\\
\Cl_3(\pi/4)=-\frac3{256}\zeta(3)+\frac1{\sqrt2}L(3,\chi_8),\\
\Cl_3(\pi/6)=\frac1{24}\zeta(3)+\frac{\sqrt3}{2}L(3,\chi_{12}).
\end{gather*}
Here $\chi_D$ is the primitive even quadratic character of conductor $D$.
\end{proposition}
These are exact coordinate identities from the residue-class Fourier sum.
No theorem asserting that the three $L$-values are independent of
$\zeta(3)$ or $\pi^3$ is supplied. The source searches at height $10^{10}$
are numerical evidence in a stated basis, not an ``if and only if''
classification of all possible reductions.

\section{The fundamental constants, by conductor}
The residue classes of a cyclotomic series can be organized either by
additive Fourier modes or by Dirichlet characters. Odd Clausen values use
the odd character sector. Distribution relations first remove non-coprime
residues; the remaining finite character transform has $\varphi(n)/2$
formal odd coordinates. Its coefficients are generally algebraic, so this
coordinate count is not a proof of the dimension of the numerical
$\mathbb Q$-span of the periods.

For the levels $8,12,24$, convenient generators are
\[
V=\Cl_2(\pi/3),\quad G=\Cl_2(\pi/2),\quad
W_8=\Cl_2(\pi/4),\quad W_{24}=\Cl_2(\pi/12).
\]
There are no primitive odd characters of conductor $12$; its odd characters
are induced from conductors $3$ and $4$. This explains the distribution
reductions below. Conductors $3$ and $4$ each have one primitive odd
character, while conductor $5$ has two conjugate primitive odd characters.
The source's printed sequence of primitive-character counts started with
two incorrect zeros; it is not used here.

Their canonical $L$-value (trigamma) forms are
\begin{equation}\label{clausen:eq:trig}
  G=\frac{\psi^{(1)}(\tfrac14)-\pi^2}{8},\qquad
  \sqrt3\,V=\frac{3\psi^{(1)}(\tfrac13)-2\pi^2}{6}\ \ \Bigl(\text{i.e. }\psi^{(1)}(\tfrac13)=\tfrac{2\pi^2}3+2\sqrt3\,V\Bigr).
\end{equation}

\section{Reduction tables and the settled level-24 question}
Every $\Cl_2(2\pi a/n)$ at $n\in\{8,12,24\}$ reduces to the basis above (PSLQ, $\ge110$ digits,
Champernowne canary; all confirmed in Wolfram):
\begin{align*}
\text{level }12:\quad
  &\Cl_2(\tfrac\pi6)=\tfrac23G+\tfrac14V,\quad \Cl_2(\tfrac{5\pi}6)=\tfrac23G-\tfrac14V,\quad
  \Cl_2(\tfrac{2\pi}3)=\tfrac23V;\\
\text{level }8:\quad
  &\Cl_2(\tfrac{3\pi}4)=W_8-\tfrac12G;\\
\text{level }24:\quad
  &\Cl_2(\tfrac{5\pi}{12})=\tfrac{16W_8+24W_{24}-4G-3V}{24},\quad
  \Cl_2(\tfrac{7\pi}{12})=\tfrac{4W_8+6W_{24}-3G}{6},\\
  &\Cl_2(\tfrac{11\pi}{12})=\tfrac{24W_{24}-8G-3V}{24}.
\end{align*}
The displayed formulas prove that the level-$24$ values are generated by
$G,V,W_8,W_{24}$. They give an upper bound of four on their rational span.
The source did not prove that $W_{24}$ lies outside the span of the other
three, so the minimal number of numerical generators remains open.
This distinction is central when a symbolic conductor count is compared
with a PSLQ experiment.

\section{A new closed form}
Since $\Cl_2(\pi/12)=W_{24}$ is itself fundamental, the cleanest carrier is the dilogarithm at the
primitive $24$th root $e^{\iu\pi/12}$:
\begin{equation}\label{clausen:eq:new}
  \Li_2(e^{\iu\pi/12})=\frac{73\pi^2}{576}+\iu\,W_{24},
\end{equation}
the real part being $\tfrac{\pi^2}6-\tfrac{\pi}{2}\!\cdot\!\tfrac\pi{12}+\tfrac14(\tfrac\pi{12})^2$. Pulling
this back through $z=1/(1-e^{\iu\theta})$ gives Reshetnikov-style closed forms at the corresponding
$\mathbb Q(\sqrt2,\sqrt3,\sqrt6)$ points; the cleanest stays the root-of-unity form~\eqref{clausen:eq:new}.

\section{A weight parity swap}
The same character parity that makes $\Cl_2$ (weight $2$, odd characters) carry deep constants reverses one
weight up. For primitive odd $\chi$, the positive odd integer values are algebraic multiples of $\pi^s$;
since $\Imag\Li_3(e^{\iu\theta})=\sum\sin(k\theta)/k^3$ pairs with odd characters at $s=3$,
\begin{equation}\label{clausen:eq:swap}
  \Imag\Li_3(e^{2\pi\iu a/n})\in\mathbb Q\cdot\pi^3\quad\text{for all }a,n,
\end{equation}
e.g. $\Imag\Li_3(e^{2\pi\iu/3})=\tfrac{2\pi^3}{81}$, $\Imag\Li_3(\iu)=\tfrac{\pi^3}{32}$,
$\Imag\Li_3(e^{\iu\pi/3})=\tfrac{5\pi^3}{162}$, $\Imag\Li_3(e^{2\pi\iu/5})=\tfrac{4\pi^3}{125}$. Thus at
weight $3$ the \emph{imaginary} part at roots of unity is elementary, and the deep constants
($\zeta(3)$ and the even-character $L(3,\chi)$) live in the \emph{real} part---the exact reverse of weight
$2$. (On the line $\Real z=1$ the picture is different again: there $\Imag\Li_3(1+\iu)$ is the central-binomial
Gaussian tower constant, because $1+\iu$ is not a root of unity; see \texttt{polylog-vertical-line}.)
